% ================================================================
%  第 9 章 收获 (后半) 原书 p.197-208, 边码 95-110
%  源文: _work/src1/task-ch09b.md
%  接 part1-ch09a.tex (p.187-196)
% ================================================================
%
% ---------------- p.197 ----------------
\origpage{95} 反思的节奏固定下来,并充当注意力的坐标与集合点——否则我的注意力往往倾向于四散随风;而且,书写还给我们留下一道刚刚做完的工作的有形痕迹,我们可以随时求助于它。在一场旷日持久的冥想里,能够回头翻看那些书面痕迹常常是有用的:它们见证着前几天、甚至几年前某个时刻的冥想。

\origpage{96} 思想,以及对思想的细致表述,因此在我迄今所实践的冥想中扮演着重要角色。但冥想并不因此就限于一种单纯的思想工作。单凭思想,无力把握生命。思想特别擅长的,是侦测矛盾——在我们对自身的看法里、在我们与他人关系的看法里,那些往往大得滑稽的矛盾;但常常,它不足以把握这些矛盾的意义。对于被求知之欲所驱动的人,思想是一件往往有用而有效的工具,甚至不可或缺——只要人始终意识到它的限度,这限度在冥想中显而易见(在数学工作中则更为隐蔽)。要紧的是:在另一样东西出现的那些敏感时刻,思想要懂得退开,踮着脚尖消失——那另一样东西也许以一阵突然而深沉的情绪的形式出现,而与此同时,手也许还在纸上继续奔跑,在同一刻给它一个笨拙而结结巴巴的表达……

\section{(37) 惊叹}

这段关于发现冥想的回顾,完全是意料之外地来到这里的,几乎违背我自己的意志——我起初打算考察的完全不是这个。我本来想谈惊叹。那个夜晚如此丰富,满载着如此多的东西,也满载着面对这些东西的惊叹。其实在工作进行当中,就已经有一种惊叹:面对每一条被揭出的新遁词,一种难以置信的惊叹——那就像一件用粗白线草草缝成的蹩脚戏装,而我竟,简直难以相信!还一本正经地把它当成真得不能再真的东西!此后许多次,在随后的年月里,我又重新遇到这同一种惊叹,如同在冥想的第一夜,面对我所发现的事实的巨大,以及那些让我一直对之视而不见的托辞的粗劣。正是从它那些滑稽可笑的侧面,我开始发现我身上那个未曾被猜想到的世界——一个随着日子、月份与年岁的推移而显露出惊人丰富的世界。然而就在那第一夜,值得我惊叹的就已经不只是些闹剧式的插曲。正是在那一夜,我头一次重新接触到一个被遗忘的、沉睡在我身上的能力;它的性质我还不清楚,只知道它是一种能力,而且随时听我支配。

而此前那几个月,也已经充满了一种无声的惊叹——惊叹于我身上一直带着的某样东西,那东西想必从来就在那里,而我刚刚才重新与它接触。我感受到这东西,不是作为一种能力,而更像是一种隐秘的温柔,一种既极其安宁又令人不安的美。后来,在我发现自己那份被忽视了如此之久的能力的狂喜中,我忘掉了这几个默默孕育的月份——它们只留下几首散落的诗作证:一些情诗,它们夹在我那些冥想笔记中间,大概常常显得格格不入……

只是许多年之后,我才回想起这些惊叹于世界之美、惊叹于我感觉到安歇在自己身上的那份美的时光。那时我才知道:我在自己身上感觉到的这份温柔与美,以及我不久之后发现、并深刻改变了我一生的那份能力,乃是同一件事的两个不可分割的方面。

% ---------------- p.198 ----------------
而我现在也看出:在我们身上那多重的东西——创造性——之中,那温柔的、内敛的、沉默的一面,自发地通过惊叹表达出来。而人认识所爱的女人、她也认识他,同样是在对由所爱者揭示出来的、自身之中那份不可言说的美的惊叹之中。当对所探究之物、对所爱之人的惊叹缺席时,我们与世界的拥抱就被削去了其中最好的部分——被削去了那使它成为对自己、对世界的祝福的东西。不是惊叹的拥抱,是一种无力的拥抱,只是占有姿态的简单复制。它无力生出别的东西,只会生出更多的复制品——也许更大些、更粗些、更厚些,无所谓,永远不是一次更新\fnoteP{8}{(34)}{\emph{无力的拥抱}\par “拥抱”一词在我这里绝不是一个简单的比喻,日常语言在此成了一个深层同一性的反映。人们可以说——而且不无道理——那么“没有惊叹的拥抱是无力的”这话就不对:若按字面理解,大地早该人烟断绝、一片荒芜了。极端的例子是强奸,其中确实没有惊叹,然而有时一个存在就在被强奸的女人身上被孕育出来。由这样的拥抱生出的孩子,想必不能不带这印记,它属于他所分得的那份“包裹”,属于该由他去承担的东西;这并不妨碍一个新存在确实被孕育、并且出生了,确实有过创造,那是力量的标志。也确实,我曾见过某个充满自得的数学家发现并证明出漂亮的定理,那是一次并不缺乏力量的拥抱的标志!但同样真实的是:如果这位数学家的生命被他的自得闷死(在某个时期,我自己的生活在一定程度上就是如此),那么这些与数学的拥抱之果实,对他、对任何人都不是福泽。对一个强奸所生的孩子的父亲和母亲,也可以这么说。我若说“无力的拥抱”,首先指的是无力在那个自以为在创造的人身上生出更新:他其实只造出一个产品,一件外在于他的东西,在他自身之中没有深层的回响;这产品远不能解放他、不能在他身上造出和谐,反而把他更紧地捆在他所囚于其中的那份自得上,那自得不断推着他去生产、再生产。这是深层意义上的一种无力,藏在一种“创造性”的外表背后——那“创造性”说到底只是一种没有刹车的生产力。\par 我也有充分的机会认识到:自得、无力惊叹,其性质是一种真正的盲目,是天然的敏感与嗅觉的阻塞;这阻塞即使不是完全的、永久的,至少在某些特定情境中显而易见。这是一种状态,其中某位声名显赫的数学家有时会在恰恰是他最擅长的东西上,显得比最迟钝的小学生还要愚笨!在另一些场合,他又会做出技术精湛的奇迹。然而我怀疑他是否还有能力去发现那些能够更新一门学科或一门科学的简单而显然的东西。它们在他下面太远,他已不屑于再看它们一眼!要看见谁也不屑于看见的东西,需要一种他已经失去、或被他放逐了的无辜……这肯定不是偶然:近二十年来数学生产的惊人增长,新结果纷至沓来,令那位只想稍稍“跟上形势”的数学家不知所措;然而(就我从各处得到的回声所能判断的),我曾多少熟悉的那些重大思考主题,几乎没有一个是真正更新了的,没有一个是大幅度的转化(而不只是积累)。更新不是量的事,它与投入的量无关——那量可以用某些“水平”的某些数学家花在某个题目上的“数学家人日”数目来衡量。一百万个“数学家人日”也无力生出一件像零那样孩子气的东西,而零更新了我们对数的感知。只有无辜有这份力量,它可见的标志之一,就是惊叹……} (34)。正是在我们还是孩子、准备为世上事物之美和为我们自身之美而惊叹的时候,我们才也准备好了去更新,并准备好作为柔软而顺从的工具,握在工匠手中,好让经祂的手、并借我们,也许有些存在与事物得以更新。

我清楚地记得,在四十年代末及随后的年月里,在这个对我而言代表着数学界的、不拘礼数的小圈子里——它有时喧闹而自信,那种略显断然的口气并不罕见(不过其中从不掺进一丝自得)——在这个圈子里,任何时候都为惊叹留有一席之地。惊叹在谁身上最显眼?是迪厄多内。不论是他作报告,还是他只在下面听,每当那个关键时刻到来、一道豁口突然敞开,就看见迪厄多内喜出望外、满面放光。那是纯粹的惊叹,有传染性、不可抗拒——其中“我”的一切痕迹都消失了。此刻我提起他,才意识到:这惊叹本身就是一种力量,它在他整个人周围立即发生作用,像一种以他为源头的辐射。要说我见过哪位数学家动用一种强大而基本的“鼓励之力”,那就是他!在此之前我从未再想到这一点,但现在我记起来了:也正是带着这样的心情,他当初接纳了
% ---------------- p.199 ----------------
\origpage{97} 我在南锡的最初几项结果,解决了他和施瓦茨一起提出的问题(关于 $(F)$ 空间与 $(LF)$ 空间)。那些结果都很平常,既不天才也不出众,可以说并没有什么值得惊叹的。后来我见过气派完全不同的一些东西,被那些自以为是大人物的同行一句不容分说的轻蔑打掉。迪厄多内身上全然没有这类负担——无论那自命是否站得住。没有任何这类东西妨碍他连小事也为之欣喜。

\origpage{98} 这种欣喜的能力里有种慷慨,它既造福于愿意让它在自己身上舒展的人,也造福于他周围的人。这恩惠的施行,并不怀着讨谁欢心的意图。它单纯得像花香,像阳光的温暖。

在我认识的所有数学家中,这份“馈赠”在我眼里以最耀眼、最有感染力、也许也最起作用的方式出现的,就是迪厄多内;我说不准\fnoteP{9}{(35)}{这份“馈赠”不是任何人的特权,我们生来都有。当它在我身上似乎缺席时,那是我自己把它赶走了,而重新接纳它只取决于我。在我身上或在某某人身上,这份“馈赠”的表达方式与另一个人不同,也许不那么有感染力、不那么不可抗拒,但它并不因此就不在,我也说不出它是否就不那么起作用。} (35)。但在我乐于交往的数学界朋友中间,没有哪一个身上缺这份馈赠。它总能找到机会显露,也许方式更含蓄,但随时都在。每当我走向他们中的一位,想与他分享我刚发现的、让我着迷的东西,它就显露出来。

如果说我作为数学家的生活里也有过挫折与痛苦,那首先就是:在我爱过的某些人身上,再也找不回我曾在他们身上认识的那份慷慨、那份对事物之美——无论“小”或“大”——的敏感;仿佛那曾使他们整个存在颤动不已的东西,已经无声无息地熄灭,被一种自得闷死了——那种自得属于这样一个人:世界对他而言已不够美,不值得他屈尊为之一喜。

当然,也还有另一种痛苦:看见我昔日某位朋友以俯就或轻蔑的态度对待我今天的某位朋友。但说到底,这痛苦也是同一种封闭造成的。对一件事物之美——无论它多么卑微——敞开着的人,一旦感受到这美,就不由自主地也对构想或做成它的人怀有敬意。在一件出于人手的事物之美里,我们感受到做它的人身上那份美的反光,感受到他制作它时所倾注的爱。当我们感受到这美、这爱,我们心里就不可能有俯就或轻蔑——正如在一个我们感受到某个女人的美、以及这美所标志的她身上那份力量的时刻,我们不可能对她有俯就或轻蔑。

\section{(38) 回归的冲动与更新}

迪厄多内身上时或闪耀的那种欣喜,想必触动了我身上某种深沉而有力的东西,这记忆此刻才如此强烈、如此鲜活地回到我心头,仿佛我刚刚还亲眼目睹。(其实我已经将近十五年没怎么有机会见到迪厄多内了,只有一两次匆匆照面。)当然,在意识层面,我并未对此给予特别注意——那不过是我这位年长的同行兼朋友那外放的性格中一个有点动人的特点,有时几乎滑稽。相反,对我重要的是:我在他身上找到了那个完美的、我甚至可以说是梦想中的合作者,来把那些将要为我眼前展开的广阔远景充当基础的东西,以一丝不苟的细心、一种充满爱意的细心,白纸黑字地写下来。直到此刻,当我把这两件事一并提起,
% ---------------- p.200 ----------------
\origpage{99} 那条联系才突然在我眼前显现:使迪厄多内成为一项伟大任务之理想仆人的——无论是在布尔巴基内部,还是在我们为另一项奠基性大工程而进行的合作中——是慷慨,是他的工作中、在他那些重大投入的抉择中,全无一丝虚荣的痕迹。我经常看见他隐没在他甘为仆役的那些任务后面,毫不吝惜地把无穷的精力献给它们,不求任何回报。毫无疑问,他什么也不求,却在工作本身、在他投入其中的那份慷慨本身里,得到一种充实与舒展——所有认识他的人都必定感受过。

我常常感到从他身上闪耀出来的那种发现之欣喜,在我心里立刻与一种相似的欣喜联系在一起——我曾在一个很小的孩子身上目睹过那种欣喜。两个记忆争相涌来——两个都让我重新看见我那时还很小的女儿。第一幅画面里,她大概才几个月,想必刚刚开始会爬。她大概是从人家把她放坐的那块草坪,爬到了旁边一条砾石小径上。她在发现那些小石子,一种无声的——却是起作用的——狂喜:她双手大把抓起石子,往嘴里送!——另一幅画面里,她大概一两岁,有人刚往金鱼缸里撒了鱼食。鱼儿们争先恐后地游过去,大张着嘴,吞下那些悬浮着、慢慢沉向缸底的小小黄色碎屑。小家伙此前从没意识到鱼也像我们一样吃东西。她身上仿佛突然一亮,迸出一声纯粹欣喜的叫喊:“妈妈看,它们在吃饭!”这确实值得惊叹——她在一瞬间发现了一个大奥秘:我们与一切其他生物同属一族的奥秘……

\origpage{100} 一个小孩子的欣喜里有一种无法用言语传达的、有传染性的力量,一种从他身上辐射出来、作用于我们的力量——而我们大多数时候总是竭力躲开它。在内心里安静的时刻,人能感觉到这力量随时都在孩子身上。只是有些时候它的作用更强些。这种围绕孩子的“力场”,在新生儿身上、在生命最初的日子和月份里最强。大多数情况下,它在整个童年期间都还感觉得到,随年岁渐长而一丝丝磨损,直到青春期——那时往往已经看不出痕迹了。然而,人们也能发现它辐射在各个年龄的人周围:在某些人身上,在得天独厚的时刻;在少数另一些人身上,则像一种气息或光晕,时时刻刻环绕着他们这个人。我有幸在童年认识这样一个人,一位男子,如今已经去世。

我也想到另一种力量,或能力:人们有时感到它从某个女人身上辐射出来,尤其是在她舒展于自己的身体、与身体交融的时刻。我常想到的词是“美”,它道出了其中一个侧面。这种美与任何美的规范、与所谓“完美”毫无关系;它不是青春的特权,也不是成熟的特权。它更像是一个人身上某种深层和谐的标志。这种和谐常常仍是残缺的,然而它正是通过这种辐射显现出来——那是力量的标志。这是一种把我们引向它所自出的中心的力量——或者不如说,它在我们身上唤起一种深沉的冲动:回到那位“女人—母亲”的身体里去——我们在生命破晓时,正是从那里出来的。它的作用有时强不可挡,当它来自所爱的女人时,令人翻覆。但对于不刻意对它封闭的人,它在每一个让这种美、这种深层和谐在自己身上舒展的女人身上都感觉得到。

从孩子身上辐射出来的力量,与从那个爱着自己身体的女人身上散发的力量是近亲。后者不断生出前者,正如孩子不断从母亲生出。但童年的力量,其本性不是吸引,也不是排斥。这种力量对不躲开它的人所施的谦卑而谨慎的作用\origpage{101},
% ---------------- p.201 ----------------
是一种更新的作用。

\section{(39) 夜之美、日之美(或:奥吉亚斯的牛圈)}

对我某个孩子的惊叹的记忆,是在五十年代末、六十年代初。如果说后来出生的其他孩子没有给我留下类似的记忆,那也许是因为我自己惊叹的能力已经钝了,我变得太疏远,无法与某个孩子的欣喜相通,甚至无法仅仅做一个见证者。

我还从未想过要追索这种能力在我一生中——从童年到今天——的变迁。那里想必有一条线索,一个极其灵敏的“探测器”。我从未想过追随这条线索,想必是因为这种能力的性质如此谦卑、外观几乎如此微不足道,以至于我几乎不会想到去特别注意它——我当时正全神贯注于发现和探测我所谓的我一生的“伟大力量”(它们今天仍在我身上显现)。然而,这种外观如此谦卑的能力,却提供了一个胜过一切的标志,指示着我们身上那份最稀有、最珍贵的力量之在与不在……

在我整个成年生活中,我从未与这种力量完全隔绝。无论我的生活在别的方面变得多么荒芜,我总在爱中重新找到孩子的惊叹、发现的欣喜。穿过许多片沙漠,爱的激情始终是与某个我已离开的东西之间那条活生生而茁壮的纽带,一根脐带,默默继续以温热而慷慨的血液滋养我。而在很长一段时间里,对所爱女人的惊叹,也与对她所生育的新生存在的惊叹不可分割——那些崭新的存在,无限娇嫩而生机勃勃,见证着她的力量,也承继着它。

但我这里的用意,主要是想稍稍追索这种“无辜之力”在我作为数学家的生活中的变迁,即我属于“数学家的世界”那个时期, 1948到1970年。当然,惊叹从未像在爱的激情中那样渗透我的数学激情。奇怪的是,若要我回想数学工作中某个欣喜或惊叹的特定时刻,我一个也想不起来!我接近数学的方式,从十七岁全力投入时起,就是给自己提出大任务。这些任务从一开始就总是“整理”、大扫除的任务。我看见一片表面的混乱,一堆杂乱异质之物,有时是不可称量的雾气,它们显然应当有一个共同的本质,藏着一个秩序、一种尚被遮蔽的和谐,要靠耐心、细致、往往旷日持久的工作把它揭示出来。这工作常常先是拖把和大刷子的活计——那粗活已经耗去大量精力——然后才轮到用掸子做收尾;收尾我不那么着迷,但也自有其魅力,而且无论如何显然有用。日复一日的工作中有一种强烈的满足:看见那个被猜到的秩序一点点显露出来,它总是比先前瞥见和猜到的更为精致,质地更为丰富。工作不断充满意外插曲,它们多半从考察某个看似极微小、此前一直被忽视的细节中冒出来。常常,对某个“细节”的打磨,给多年前做过的工作投下意想不到的光。有时它也带来新的直觉,而对它的深化又成了另一项“大任务”的对象。

这样,在我的数学工作中(除了1954年前后那个我有机会谈过的“艰难的一年”),有一种持续的悬念,注意力始终被吊着。何况,对我那些“任务”的忠诚不许我跑得太远,我咬着嚼子,急不可耐地想
% ---------------- p.202 ----------------
\origpage{102} 把它们全都做完,终于纵身跃入未知、真正的未知——而这些任务的规模此时已经变得如此之大:要把它们做完,即便有那些终于赶来救援的善意相助,我的余生也不够用!

我的主要向导,是不断寻求一种完美的融贯、一种完整的和谐——我猜到它就在事物动荡的表面之下,我耐心地努力把它揭示出来,从不厌倦。那确实是一种敏锐的“美”感,它是我的嗅觉,也是我唯一的罗盘。我最大的快乐,与其说是在它出现在光天化日之下时凝视它,不如说是看着它一点点从它喜欢不断躲藏其中的阴影与雾气的衣袍下显露出来。当然,不把它带到最明亮的日光下,我是不肯罢休的。那时我有时也体会到凝视的圆满——当一切可闻的声音汇成同一片浩瀚的和谐。但更常见的是:被带到日光下的东西,立刻又成了再一次潜入雾中的动机与手段,去追寻那位永远神秘、永远未知的“她”的又一次化身——她不断呼唤我,为了再认识她……

\origpage{103} 迪厄多内的快乐与欣喜,我觉得,主要在于看见事物之美在光天化日之下显现;而我的快乐,首先在于在雾与夜的皱褶里追寻它。这也许就是迪厄多内与我接近数学的方式之间的深刻差异。对事物之美的感觉,至少在很长一段时间里,在我身上想必不比他弱——尽管它也许在六十年代里、在一种自得的作用下变钝了。但看来,在迪厄多内身上通过惊叹显现出来的对美的感知,在我身上采取了不同的形式:不那么静观,更爱动手,在所感受和表达的情绪层面上也不那么外显。若果真如此,我要追索的,就应当是我身上对数学事物之美的那份敞开之变迁,而不是那份神秘的“惊叹之馈赠”。

\section{(40) 竞技数学}

相当清楚的是,对数学事物之美的敞开,在我身上从未完全消失,即使在六十年代直到1970年——那时自得在我与数学、与其他数学家的关系中逐步占据了越来越大的位置。若没有最低限度对事物之美的敞开,我根本不可能作为数学家“运转”起来,哪怕只开到最低档——而且我怀疑,任何人若身上不多少还活着这种美感,能否在数学里做出有用的工作。依我看,造成这位数学家与那位数学家之间、或同一位数学家的这项工作与那项工作之间差别的,与其说是某种所谓的“脑力”,不如说是这种敞开或敏感的精细度、细腻度——研究者彼此之间不同,同一个研究者不同时刻也不同。最深刻、最丰产的工作,也是那种见证着最敏锐的敏感性的工作——那敏感性用于把握事物隐藏的美\fnoteP{10}{(36)}{在我看来,这种对美的精细敏感,与我曾有机会以“严苛”(对自身)或“严格”(就其完整意义而言)之名谈过的东西密切相连;我曾把它描述为“对我们自身之中某种精细之物的注意”,一种对所受探究之物的理解品质的注意。对一件数学之物的这种理解品质,不能与对该物特有之“美”或深或浅、或完美或粗略的感知分开。} (36)。

若果真如此,那就得相信,这种敏感在我身上一直鲜活到最后,至少某些时刻如此——因为正是在六十年代末\fnoteP{11}{}{(8月8日)核实之后看来,我关于 Motif 的思考之开端是在六十年代初,而不是六十年代末。},我开始瞥见并稍稍揭示出我所发现的最隐蔽、最神秘的数学对象:我称之为“Motif”的东西。
% ---------------- p.203 ----------------
\origpage{104} 它也正是在我作为数学家的生活中对我施以最大魅惑的东西(若不算最近几年的一些思考——它们其实与 Motif 的实在密切相关)。毫无疑问,若不是我的生活突然转向一条完全意外的道路,把我远远带离数学事物那宁静的世界,我最终会追随这强大魅惑的召唤,把那些一直把我囚禁着的“任务”撂在一边!

也许我可以说:在我工作室的孤独里,美感一直保持不变,直到1970年我第一次“觉醒”的时刻,并没有真正受到那种常常标示着我与同类关系的自得的影响?某种“嗅觉”甚至想必随年岁而愈发精细,在与数学事物每日而亲密的接触中。我们对事物所能拥有的亲密认识,有时使我们能把握超出当下所知的东西,并更深入地进入认识——这种认识或这种成熟,以及作为其最显眼标志的这种“嗅觉”,是向事物之美与真之敞开的近亲。它助长、刺激这样的敞开,而且它是此前一切敞开时刻、一切“真理时刻”的总和与果实。

那么,我接下来要考察的,就是:这种对美的自发敏感,在我与这位或那位同行的关系中、在它有机会显现的那些时刻,在多大程度上受到了或深或浅的扰乱。

记忆在这方面交给我的东西,并不能凝成一个有形的、确切的事实,让我在这里多少详尽地摆出来。这里的记忆仍然只是一团雾,但它还是交给我一个总体印象,我得设法把它界定出来。那是某种内心态度留在我身上的印象,它想必最终变成了第二天性,每当我接到一条数学信息、而所涉之物多少“在我的射程内”时,它就显现出来。说实话,从某个相对无害的侧面看,这种态度想必一直属于我,它属于某种气质,我也曾有机会顺带触及。说的就是这种反射:首先只肯了解一个陈述,绝不看它的证明,先试着把它安放到我所已知的东西里,看看用这些已知的东西来说,这陈述是否变得透明、显然。这常常使我以或深或浅的方式重新表述那个陈述,使之更一般或更精确,往往两者兼有。只有当我无法用我的经验和我的意象去“安放”这个陈述时,我才准备好(有时几乎违背我自己的意愿!)去听(或读……)那些来龙去脉——它们有时给出事情的“那个”道理,或者至少给出一个证明,懂不懂另说。

这是我接近数学的一个特点,我觉得它使我有别于当时布尔巴基的所有其他成员,也使我实际上不可能像他们那样插入集体工作。这特点想必也在我教书时构成一种障碍,我的所有学生想必都感受过——直到今天(多亏上了年纪)它才终于松动了一些。

我身上这个特点,想必已经是一种敞开上的欠缺。它只意味着一种部分的敞开,只肯接纳“来得正好”的东西,至少对别的一切都极不情愿接纳。在我的数学投入之抉择中,在我同意把时间花在这样或那样不期而至的信息上的抉择中,这种有意“部分封闭”的打算,今天比任何时候都更强。它甚至是一种必要,如果我想能追随那最令我着迷之物的召唤,而又不再把“我的命拿去吃”献给数学夫人!

% ---------------- p.204 ----------------
\origpage{105} 然而这团“雾”交给我的,不止这个特点——我几年前就已经察觉它了(晚知道总比不知道好!)。在某个时刻,这种反射变成了某种关乎荣誉的事:我要是在比说这句话更短的时间里还“拿不下”这个陈述(假定它对我还不算太熟悉),那才叫见鬼!如果陈述的作者是个有名有姓的陌生人,还多一层微妙:那可不行,我(毕竟人家以为我是在行的人!)岂能手里没有这一套!而常常我确实有,而且绰绰有余——那时我的态度就倾向于这个方向:“行了,您可以回去穿好衣服了——等您做得好一点再来吧!”

这恰恰就是我在“那个在我园子里乱走的毛头小子”一事上的态度。我甚至不敢发誓:他做的东西里没有一些有趣的细节,是我那些“秘密笔记”里所做的东西不曾覆盖的——不过这终归是无关紧要的事\fnoteP{12}{}{(8月8日)后来我发觉,这件事并不那么“无关紧要”:它构成了从“竞技态度”通向某种不诚实之开端的界线,而这条界线,我也许有几次越过了……}。说到底,这段插曲也照亮了我在这里考察的问题:对数学事物之美的那份敞开之深层扰乱。仿佛从我把某样东西“做出来”的那一刻起,它的美对我而言就消失了,剩下的只有一种虚荣,要求为它记功受利。(而我却又不屑于花时间去发表它——诚然,要发表的也太多了。)那是一种典型的占有态度,类似一个男人,认识了一个女人之后,再也感觉不到她的美,又去追逐另外一百个,却仍容不得别人认识她。这种态度我在爱情生活中是斥责的,自以为远在这样一份虚荣之上,却又小心翼翼地不去确认这个显而易见的事实:我对数学的态度恰恰正是如此!

\origpage{106} 我有个印象:我刚才在自己身上指出的这些粗鄙的竞争心性、不妨说“竞技”心性,想必大约就在它们于我自己身上变得普遍的那个时候,开始在我的数学“圈子”里变得普遍起来。要我在时间上定位它们出现的时刻,或它们变成这个圈子里所呼吸的空气之一部分的时刻,或我的学生们在与我接触中所呼吸的空气的时刻,我却无能为力。我唯一敢说的,是这大概在六十年代,也许早在六十年代初,或五十年代末。(若果真如此,我的学生全都摊上了——他们只能要么接受,要么拉倒!)要定位它,我还需要别的确切例子,而此刻它们完全逃出了我的记忆。

当然,这个卑微的实情,与我为自己同数学、同一般年轻研究者的关系所描绘的高贵形象完全相反。用来蒙骗我自己的那个粗劣托辞,灵感来自功绩主义:对于这个形象,我所留住的只是我与学生们的关系(他们为我的声望添砖加瓦,是我声望中最高贵的花朵!),以及我与那些特别出色的年轻数学家的关系——我懂得承认他们的功绩,并且像对待我的学生一样平等待他们,不等他们的头上戴上桂冠(当然,桂冠也没等多久——嗅觉这东西,有就是有,没有就是没有!)。至于那些既没福气做我学生、也没福气做我某位朋友的学生、又不是年轻天才的年轻人,我丝毫不在意我与他们是什么关系。他们不算数。

我相信,当我与年轻研究者有了个人的接触——无论是在讨论班上遇见他,还是他写信给我——上述实情多半会松动、缓和一些。从这一点看,“毛头小子”那件事也许是个稍显特殊、例外的情形。我觉得,对我刚才谈到的那些研究者,我大概会多少把他们看作“受我庇护”的人,这想必在我身上唤醒了一种更和善的态度。在这种情形里,我想出头的欲望也能找到出口:对当事人发表我的评论,建议他也许从更广阔的视角、或者更深地进入事情本身去重做他的工作。在这种情形下,那位在有限时间里多少扮演了学生角色的年轻研究者,很可能也从中有所得,并且对我与他的关系留下好印象。(这方面任何或正或反的回声,我都欢迎。)

% ---------------- p.205 ----------------
\origpage{107} 我在这里主要想到的是较年轻的研究者,而“竞技”态度绝不限于我与他们的关系,这不用说。但确实是在与年轻研究者的关系中,一位有名望的数学家的影响——无论心理上还是实际上的——往往最强,对他们未来的职业生涯也最有后果。

\section{(41) 克里希那穆提,或:成为羁绊的解放}

\origpage{108} 这一夜,我停在一份宽慰之情、一份极大的满足上——那是一个没有白费时间的人的满足!我突然觉得轻松而快乐——一种有时带点促狭的快乐,迸成顽皮的笑——一个捣蛋鬼恶作剧的笑。其实我并没做什么大事,我不过是从一个略微不同的角度,重看了一段已经“知道”的插曲,即那个著名的“毛头小子……”。这个角度显示出的,是我与数学本身的关系——在某些情形下,不只是我与数学家的关系。就这一点,已经足够让一个我曾珍爱的神话化为青烟。

说实话,这并非我第一次审视我与数学的关系。两年半以前,我就曾被引向为此花上几个星期或几个月。那时我(在许多别的事情之外)意识到:在我过去对数学的投入中,自我之力、自我膨胀之力有多么重要。但昨夜我刚刚点到了一个当时逃过我注意的侧面。现在我回过头来谈它,发现这个侧面——即我与数学的关系中那种嫉妒态度——与第一夜“冥想”结束时到来的那个“再平常不过”的发现接上了(当时我并不知道自己在冥想,就像茹尔丹先生不知道自己是在说散文……)。这很可能在我随后那份欢欣的狂喜中占了一份。即便没有被自觉地察觉,那也多少像是对我从前发现的某样东西在新光照下的再次确认——而这时心里的快乐,与数学里的一样:当你并没有去找,却从一条完全不同的路径撞上你认识的、也许多年前就发现的东西。每一次,这都伴随着一种深切的满足感:事物的和谐再一次显露,而我们对它的认识也或多或少被更新。

而且我相信,这一次我确实“转完了”!好几天来我都清楚地感到还有什么东西没被挖出来,却又说不清到底是什么。我没有试着硬来,我觉得只需让它自己来,让所追随的那条线自由展开,穿过一片片既熟悉又意外的风景。说意外,是因为我此前从未费心去看它们。我是以散步的步子走近那个剩下的“热点”的。而我相信,这是我刚刚走完、且已接近尾声的这趟旅程中的最后一个。

而刚一到达这一点,我就有了一种感觉,像一个登上观景台的人:从那里,他看见\origpage{109}自己刚刚走过的那片风景铺展开来,而在途中每一刻,他只能看见其中一角。
% ---------------- p.206 ----------------
\origpage{110} 而现在有了这种辽阔与空间的感觉,那是一种解放。

如果我试着用文字表述眼前这片风景交给我的东西,出来的就是这些:这些年来,在我作为数学家的生活中到来的一切——常常来得不是时候,也不受欢迎——都是我在还属于数学家的世界那些年月里所播种之物的收获与讯息。

当然,这件事,这些年来我已经对自己说过又说,在我刚写下的这些笔记里也说过。我对自己说过,多少类比于另一些执拗地到来的收获——我长久拒斥,最终接纳并化为己有。从我这样接纳的第一份收获起,甚至在我懂得冥想之前,我就明白了:每一份收获都必有其意义,而赌气只是回避一个意义、推迟一个了结的期限。这一认识对我很宝贵,因为它常使我免于自怜,也免于那种义愤——义愤往往是自怜的一种伪装形式。这认识在我身上像是一种半成熟,它还远没有终结那种积习难改的反射:当收获显得苦涩时加以拒绝。当我说“赌气没有用”,收获并没有因此被接纳。我也许既不自怜也不义愤,可我仍在“赌气”!只要那盘菜没吃下去,它就还没被接纳——而不吃,就是赌气。

接纳并吃下去是一件工作:某种能量在“工作”,一种工作在光天化日下或在暗处进行,有东西在转变……而赌气则是一种能量的浪费,那能量散掉了——散在“赌气”上!而吃、消化、吸收这件工作,是省不掉的。仅仅经历一些事件、仅仅“取得”或“获得”一种经验,与一件工作毫无共同之处。它只是一份可供工作的材料,做不做由你。从我遇到数学家的世界起,三十六年了,我一直在使用这份自由:回避一件工作,而可供吃、可供消化的材料却年复一年地增多。我从昨天起感到的这份欢欣的解放感,确实是一个标志:我面前那件工作——我不断地以别的工作或任务为借口推开的——终于做完了。也确实到时候了!

现在还不能确定事情果真如此,还不能确定没有剩下某个阴暗而顽固的角落逃过了我的注意、我还得回到那里去。但同样真实的是:这种解放感不会骗人——我一生中每一次感到它,事后都能证实它确实是一次解放的标志;是某种持久的、已经获得的东西,是一种理解、一种认识的果实,而那认识已成为我的一部分。若我愿意,我有自由无视这认识,把它埋在我想要的地方、按我想要的方式埋掉。但要毁掉它,却不在我、也不在任何人的能力之内——正如人不能毁掉一只果子的成熟,不能把它弄回一种已不属于它的青涩状态。

再一次证实我并不比别人“更好”,这是一大宽慰。当然,这也是我相当经常地对自己重复的话——可是重复和看见,毕竟不是一回事!没有孩子的无辜与灵动——他看见,就像他呼吸——要看见明显之事,往往需要一件工作;好了,做完了,我终于看见了这一件:我并不比几天前还让我“喘不过气”的那些同行或旧学生“更好”!想想我如今卸下的这副担子有多重吧!自认为比别人强,从某种意义说也许挺受用,但那也非常累人。那甚至是精力的极大浪费——正如每一次要维持一个虚构时那样。人很少意识到这一点,但光是顶着风浪维持那个虚构,就已经需要精力了——而明显之事在每一场合都冲着我那仔细塞住的耳朵喊:这是假的,看看吧,傻瓜!看见,有时也许是一件工作,
% ---------------- p.207 ----------------
但一旦做完就做完了。这就一劳永逸地省掉了两件事:一是像那样动不动就堵上眼睛和耳朵地四处走动——这本身也是桩苦差!二是每当有件我自己不小心搁在那里的东西落到我头上时,就把它当作不可容忍的侮辱而自苦。

受够了这套把戏!看清了把戏,就已经出来了。付了钱,没错,我有权一直在里面转到老死,甚至有义务——那也无妨,所有人都会这么对我说:权利、义务——看人下菜碟。所有这些是权利又是义务、是义务又是权利的东西,也累人得很,当我自认为比别人强的时候,它们就黏上我。说到底这很正常:既然是强者,就悄悄收着(这就是“权利”),并且“付账”,为了人类精神与数学的荣誉尽全部义务——真的很美,没错,荣誉、精神、数学,谁还能说得更好,好极了!再来一个!是很美,没错,可也非常累人,到头来会让人落枕。我落过枕了,现在够了——我把位置让给别人去挺直腰板吧。

这也正常(既然我谈到学生):学生超过老师。我曾为此恼怒,我真有精力可浪费!这些都结束了!

多么轻松!

% ---------------- p.208 ----------------
